\documentclass[11pt]{article}
\usepackage[margin=1in]{geometry}
\usepackage[T1]{fontenc}
\usepackage{hyperref}
\title{Tactix interim audit: online configuration boundary}
\author{Codex for Andrew}
\date{25 September 2026}
\begin{document}
\maketitle

\paragraph{Outcome.} The codebase is not yet ready for production online
learning. This source-only checkpoint changed no solver or acceptance contract.
The detailed evidence is in the AM438 technical report under
\texttt{benchmarks/reports/}.

\paragraph{Configuration lifetime.} All three public online-configuration
entry points finalize C input before rejecting an invalid program. The online
executor's own replacement is transactional, but the surrounding input
session can become closed: later checked clause additions then return -1.
The C header does not promise rollback on rejection, so this is a contract
risk, not a claimed regression. A simple call reorder is unsafe because input
finalization retires absent variables and settles proof identities used by
program conversion. A design should specify admission versus finalization and
cover valid and invalid v1--v3 programs with attached and unattached proofs.

\paragraph{Learning architecture.} A synchronous solve currently creates a
fresh stack-local DefaultPolicy. The persistent C session owns a separate
online inprocessing executor, not the running default policy. For in-run
adjustment of pass dose/order and switching branching or clause-deletion
sub-strategies, the next structural work is a session-owned policy with
default-exact safe-point slots and owned strategy arms. The current reduction
seam chooses victims only; collection and tier decisions remain embedded in
DefaultPolicy. Stage this behind a byte-identical default trace gate, then
check proofs, models, public reachability, fuzz, and matched performance at
each cutover. This is a recommendation, not an implemented improvement.

\paragraph{Correctness gates.} Current-source Debug job 1723391 had passed
186 of 187 tests at this snapshot and was still running the final scripted
demo. Both AM434 extended AES proof-checking tasks were still running. Their
terminal results must be inspected; no acceptance follows from partial
passes. The reduced H2 native DSR still fails the canonical no-option checker,
so N2 remains open. N1, N3, N4, affected fuzz, and final learning readiness
also remain open. The requested 1,166-instance Tier D, 5,000-second,
4-CPU/32-GiB hwulab comparison is correctly deferred until those goals are
actually verified.

\paragraph{Next actions.} Finish the running gates; resolve the general H2
proof/checker-state discrepancy under a chosen canonical contract; then land
the staged policy/session refactor with exact default-parity checks. Do not
promote a partial checker result or run Tier D early.
\end{document}
